\documentclass[11pt]{article}

\usepackage{amsmath,amssymb}

\usepackage[margin=24mm]{geometry}

\usepackage{longtable}

\allowdisplaybreaks

\newcommand{\Sfun}{\mathcal{S}}

\newcommand{\Cfun}{\mathcal{C}}

\newcommand{\Tfun}{\mathcal{T}}

\title{SoftArm PCC Model (1 Section)}

\author{Generated by SoftArm Toolbox}

\date{}

\begin{document}

\maketitle

\section{Model Summary}
This document is generated from the same SymPy model used for numerical code generation.
\begin{description}
\item[Source configuration] \texttt{pcc\_flying\_plane\_contact\_n1.toml}
\item[Arm family] \texttt{pcc}
\item[Number of sections] 1
\item[Base mode] \texttt{floating\_rpy}
\item[Arm inertia model] \texttt{lumped}
\item[Material integration] \texttt{analytic}
\end{description}

\section{Notation and Parameters}
\begin{equation}
\boldsymbol q_B=\begin{bmatrix}x_B\\y_B\\z_B\\\phi_B\\\theta_B\\\psi_B\end{bmatrix}
\end{equation}

\begin{equation}
\dot{\boldsymbol q}_B=\begin{bmatrix}\dot{x}_B\\\dot{y}_B\\\dot{z}_B\\\dot{\phi}_B\\\dot{\theta}_B\\\dot{\psi}_B\end{bmatrix}
\end{equation}

\begin{equation}
\boldsymbol q_a=\begin{bmatrix}b_{x,1}\\b_{y,1}\\L_{1}\end{bmatrix}
\end{equation}

\begin{equation}
\dot{\boldsymbol q}_a=\begin{bmatrix}\dot{b_{x,1}}\\\dot{b_{y,1}}\\\dot{L_{1}}\end{bmatrix}
\end{equation}

\begin{equation}
\boldsymbol q=\begin{bmatrix}\boldsymbol q_B\\\boldsymbol q_a\end{bmatrix},\qquad \dot{\boldsymbol q}=\begin{bmatrix}\dot{\boldsymbol q}_B\\\dot{\boldsymbol q}_a\end{bmatrix}
\end{equation}

The NED world frame is used; gravity acts in the positive world $z$ direction.
\subsection{Runtime Parameters}
\begin{longtable}{lll}
\hline
Symbol & Code name & Default \\
\hline
\endfirsthead
\hline Symbol & Code name & Default \\
\hline
\endhead
$L_{0,1}$ & \texttt{s1\_rest\_length} & $0.45$ \\
$m_{1}$ & \texttt{s1\_mass} & $0.18$ \\
$I_{xx,1}$ & \texttt{s1\_Ixx} & $0.0018$ \\
$I_{yy,1}$ & \texttt{s1\_Iyy} & $0.0018$ \\
$I_{zz,1}$ & \texttt{s1\_Izz} & $0.0001$ \\
$k_{b_x,1}$ & \texttt{s1\_k\_bx} & $1.1$ \\
$k_{b_y,1}$ & \texttt{s1\_k\_by} & $1.1$ \\
$k_{L,1}$ & \texttt{s1\_k\_l} & $45$ \\
$d_{b_x,1}$ & \texttt{s1\_d\_bx} & $0.05$ \\
$d_{b_y,1}$ & \texttt{s1\_d\_by} & $0.05$ \\
$d_{L,1}$ & \texttt{s1\_d\_l} & $0.1$ \\
$g$ & \texttt{gravity} & $9.81$ \\
$m_e$ & \texttt{tip\_mass} & $0.1$ \\
$I_{e,xx}$ & \texttt{tip\_Ixx} & $0.01$ \\
$I_{e,yy}$ & \texttt{tip\_Iyy} & $0.01$ \\
$I_{e,zz}$ & \texttt{tip\_Izz} & $0.01$ \\
$m_B$ & \texttt{vehicle\_mass} & $1.5$ \\
$I_{B,xx}$ & \texttt{vehicle\_Ixx} & $0.03$ \\
$I_{B,yy}$ & \texttt{vehicle\_Iyy} & $0.03$ \\
$I_{B,zz}$ & \texttt{vehicle\_Izz} & $0.05$ \\
$r_{\mathrm{t1},1}$ & \texttt{act\_t1\_s1\_radius} & $0.0275$ \\
$r_{\mathrm{t2},1}$ & \texttt{act\_t2\_s1\_radius} & $0.0275$ \\
$r_{\mathrm{t3},1}$ & \texttt{act\_t3\_s1\_radius} & $0.0275$ \\
$x_{\mathrm{tool\_offset}}$ & \texttt{constraint\_tool\_offset\_x} & $0$ \\
$y_{\mathrm{tool\_offset}}$ & \texttt{constraint\_tool\_offset\_y} & $0$ \\
$z_{\mathrm{tool\_offset}}$ & \texttt{constraint\_tool\_offset\_z} & $0$ \\
$x_{\mathrm{plane\_point}}$ & \texttt{constraint\_plane\_point\_x} & $0$ \\
$y_{\mathrm{plane\_point}}$ & \texttt{constraint\_plane\_point\_y} & $0$ \\
$z_{\mathrm{plane\_point}}$ & \texttt{constraint\_plane\_point\_z} & $0.45$ \\
$x_{\mathrm{plane\_normal}}$ & \texttt{constraint\_plane\_normal\_x} & $0$ \\
$y_{\mathrm{plane\_normal}}$ & \texttt{constraint\_plane\_normal\_y} & $0$ \\
$z_{\mathrm{plane\_normal}}$ & \texttt{constraint\_plane\_normal\_z} & $-1$ \\
$\mu$ & \texttt{constraint\_friction} & $0.3$ \\
$v_s$ & \texttt{constraint\_friction\_velocity} & $0.01$ \\
$\omega_c$ & \texttt{constraint\_stabilization\_frequency} & $20$ \\
$\zeta_c$ & \texttt{constraint\_stabilization\_ratio} & $1$ \\
\hline
\end{longtable}

\section{Kinematics}
The floating-base attitude uses the ZYX convention $R_{WB}=R_z(\psi_B)R_y(\theta_B)R_x(\phi_B)$.
\begin{equation}
H_{WB}=\begin{bmatrix}\cos{\left(\theta_B \right)} \cos{\left(\psi_B \right)} & \sin{\left(\theta_B \right)} \sin{\left(\phi_B \right)} \cos{\left(\psi_B \right)} - \sin{\left(\psi_B \right)} \cos{\left(\phi_B \right)} & \sin{\left(\theta_B \right)} \cos{\left(\phi_B \right)} \cos{\left(\psi_B \right)} + \sin{\left(\phi_B \right)} \sin{\left(\psi_B \right)} & x_B\\\sin{\left(\psi_B \right)} \cos{\left(\theta_B \right)} & \sin{\left(\theta_B \right)} \sin{\left(\phi_B \right)} \sin{\left(\psi_B \right)} + \cos{\left(\phi_B \right)} \cos{\left(\psi_B \right)} & \sin{\left(\theta_B \right)} \sin{\left(\psi_B \right)} \cos{\left(\phi_B \right)} - \sin{\left(\phi_B \right)} \cos{\left(\psi_B \right)} & y_B\\- \sin{\left(\theta_B \right)} & \sin{\left(\phi_B \right)} \cos{\left(\theta_B \right)} & \cos{\left(\theta_B \right)} \cos{\left(\phi_B \right)} & z_B\\0 & 0 & 0 & 1\end{bmatrix}\label{eq:base-transform}
\end{equation}

\begin{equation}
H_{BA}=\begin{bmatrix}1 & 0 & 0 & 0.0\\0 & 1 & 0 & 0.0\\0 & 0 & 1 & 0.0\\0 & 0 & 0 & 1\end{bmatrix}\label{eq:mount-transform}
\end{equation}

\subsection{PCC Section}
\begin{equation}
\Sfun(z)=\begin{cases}\dfrac{\sin\sqrt z}{\sqrt z},&z\ne0\\1,&z=0\end{cases},\qquad \Cfun(z)=\begin{cases}\dfrac{1-\cos\sqrt z}{z},&z\ne0\\\dfrac12,&z=0\end{cases}
\end{equation}

For one section, $b_x$ and $b_y$ are Cartesian bending-angle components and $L$ is its current length.
\begin{equation}
H_i(\xi)=\begin{bmatrix}- b_x^{2} \xi^{2} \Cfun\!\left(b_x^{2} \xi^{2} + b_y^{2} \xi^{2}\right) + 1 & - b_x b_y \xi^{2} \Cfun\!\left(b_x^{2} \xi^{2} + b_y^{2} \xi^{2}\right) & b_x \xi \Sfun\!\left(b_x^{2} \xi^{2} + b_y^{2} \xi^{2}\right) & L b_x \xi^{2} \Cfun\!\left(b_x^{2} \xi^{2} + b_y^{2} \xi^{2}\right)\\- b_x b_y \xi^{2} \Cfun\!\left(b_x^{2} \xi^{2} + b_y^{2} \xi^{2}\right) & - b_y^{2} \xi^{2} \Cfun\!\left(b_x^{2} \xi^{2} + b_y^{2} \xi^{2}\right) + 1 & b_y \xi \Sfun\!\left(b_x^{2} \xi^{2} + b_y^{2} \xi^{2}\right) & L b_y \xi^{2} \Cfun\!\left(b_x^{2} \xi^{2} + b_y^{2} \xi^{2}\right)\\- b_x \xi \Sfun\!\left(b_x^{2} \xi^{2} + b_y^{2} \xi^{2}\right) & - b_y \xi \Sfun\!\left(b_x^{2} \xi^{2} + b_y^{2} \xi^{2}\right) & - \left(b_x^{2} \xi^{2} + b_y^{2} \xi^{2}\right) \Cfun\!\left(b_x^{2} \xi^{2} + b_y^{2} \xi^{2}\right) + 1 & L \xi \Sfun\!\left(b_x^{2} \xi^{2} + b_y^{2} \xi^{2}\right)\\0 & 0 & 0 & 1\end{bmatrix}\label{eq:local-transform}
\end{equation}

\begin{equation}
H_{Wi}(\xi)=H_{WB}H_{BA}\left(\prod_{k=1}^{i-1}H_k(1)\right)H_i(\xi)
\end{equation}

\begin{equation}
J_{v,i}=\frac{\partial r_{Wi}}{\partial\boldsymbol q},\qquad J_{\omega,i}^{(:,j)}=\operatorname{vex}\!\left(\operatorname{skew}\!\left(\frac{\partial R_{Wi}}{\partial q_j}R_{Wi}^{T}\right)\right)
\end{equation}

\begin{equation}
J_e=\begin{bmatrix}J_{v,e}\\J_{\omega,e}\end{bmatrix}
\end{equation}


\section{Energy and Dynamics}
For the lumped inertia option, each section contribution is evaluated at $\xi=1/2$.
\begin{equation}
M_i=\int_0^1\!\left(m_iJ_{v,i}^{T}J_{v,i}+J_{\omega,i}^{T}R_{Wi}I_iR_{Wi}^{T}J_{\omega,i}\right)d\xi
\end{equation}

\begin{equation}
M_e=m_eJ_{v,e}^{T}J_{v,e}+J_{\omega,e}^{T}R_eI_eR_e^TJ_{\omega,e}
\end{equation}

\begin{equation}
M_B=J_B^T\operatorname{diag}\!\left(m_BI_3,R_{WB}I_BR_{WB}^T\right)J_B
\end{equation}

\begin{equation}
M=M_B+\sum_{i=1}^{N}M_i+M_e\label{eq:mass-assembly}
\end{equation}

\begin{equation}
T=\frac12\dot{\boldsymbol q}^{T}M(\boldsymbol q)\dot{\boldsymbol q}
\end{equation}

\begin{equation}
V=V_g+V_{\mathrm{elastic}},\qquad V_g=-g\!\left(m_Bz_B+\sum_{i=1}^{N}m_i z_i(1/2)+m_ez_e\right)
\end{equation}

\begin{equation}
V_{\mathrm{elastic}}=\frac12\sum_{i=1}^{N}\left(k_{b_x,i}b_{x,i}^{2}+k_{b_y,i}b_{y,i}^{2}+k_{L,i}(L_i-L_{0,i})^2\right)
\end{equation}

\begin{equation}
D=\operatorname{diag}\!\left(\begin{bmatrix}0\\0\\0\\0\\0\\0\\d_{b_x,1}\\d_{b_y,1}\\d_{L,1}\end{bmatrix}\right)
\end{equation}

\begin{equation}
c(\boldsymbol q,\dot{\boldsymbol q})=\frac{\partial(M\dot{\boldsymbol q})}{\partial\boldsymbol q}\dot{\boldsymbol q}-\frac12\left(\frac{\partial(\dot{\boldsymbol q}^{T}M\dot{\boldsymbol q})}{\partial\boldsymbol q}\right)^T
\end{equation}

\begin{equation}
h=c+\nabla_{\boldsymbol q}V+D\dot{\boldsymbol q}
\end{equation}

The vehicle wrench $w_B$ is expressed in the vehicle body frame, whereas the end wrench $w_e$ is expressed in the NED world frame.
\begin{equation}
Q=S_a\tau_a+B_v(\boldsymbol q)w_B+J_e^Tw_e,\qquad B_v=J_B^T\operatorname{diag}(R_{WB},R_{WB})
\end{equation}

\begin{equation}
M\ddot{\boldsymbol q}+h=Q\label{eq:dynamics}
\end{equation}

\begin{equation}
\dot x=\begin{bmatrix}\dot{\boldsymbol q}\\M^{-1}(Q-h)\end{bmatrix},\qquad x=\begin{bmatrix}\boldsymbol q\\\dot{\boldsymbol q}\end{bmatrix}
\end{equation}


\section{Actuation}
The configured actuator family is \texttt{tendon} with 3 ordered channels.
\begin{center}
\begin{tabular}{ll}
\hline
Channel & Type \\
\hline
\texttt{t1} & \texttt{unilateral} \\
\texttt{t2} & \texttt{unilateral} \\
\texttt{t3} & \texttt{unilateral} \\
\hline
\end{tabular}
\end{center}
\begin{equation}
y(\boldsymbol q_a)=\begin{bmatrix}- r_{\mathrm{t1},1} b_{x,1} + L_{1}\\- r_{\mathrm{t2},1} \left(- 0.49999999999999983 b_{x,1} + 0.86602540378443874 b_{y,1}\right) + L_{1}\\- r_{\mathrm{t3},1} \left(- 0.50000000000000042 b_{x,1} - 0.8660254037844384 b_{y,1}\right) + L_{1}\end{bmatrix}
\end{equation}

\begin{equation}
J_a=\frac{\partial y}{\partial\boldsymbol q_a},\qquad \gamma_a=\dot J_a\dot{\boldsymbol q}_a
\end{equation}

The actuator coordinates are tendon lengths, and positive $T$ denotes tendon tension.
\begin{equation}
\tau_a=-J_a^TT
\end{equation}


\section{Active Acceleration Constraint}
The configured constraint family is \texttt{plane\_point\_contact} with 1 channels.
\begin{center}
\begin{tabular}{ll}
\hline
Channel & Type \\
\hline
\texttt{normal\_contact} & \texttt{unilateral} \\
\hline
\end{tabular}
\end{center}
\begin{equation}
A=\frac{\partial\phi}{\partial\boldsymbol q},\qquad \gamma=\dot A\dot{\boldsymbol q},\qquad Q_c=G(\boldsymbol q,\dot{\boldsymbol q})\lambda
\end{equation}

\begin{equation}
\begin{bmatrix}M&-G\\A&0\end{bmatrix}\begin{bmatrix}\ddot{\boldsymbol q}\\\lambda\end{bmatrix}=\begin{bmatrix}Q-h\\a_c-\gamma-2\zeta_c\omega_cA\dot{\boldsymbol q}-\omega_c^2\phi\end{bmatrix}
\end{equation}

For the plane-point implementation, $n$ is the unit normal directed into the free half-space.
\begin{equation}
\phi=n^T(r_c-r_0)\ge0,\qquad A=n^TJ_c
\end{equation}

\begin{equation}
v_t=(I-nn^T)J_c\dot{\boldsymbol q},\qquad f_t=-\mu\lambda\frac{v_t}{\sqrt{v_t^Tv_t+v_s^2}}
\end{equation}

\begin{equation}
G=J_c^T\left(n-\mu\frac{v_t}{\sqrt{v_t^Tv_t+v_s^2}}\right),\qquad\lambda\ge0
\end{equation}

The generated solver assumes that contact is active; a negative reaction marks the active mode as infeasible.

\end{document}
